\section{Introduction}

Anonymous credential schemes let users prove they hold certain attributes, such
as their age or nationality, without revealing who they are or anything about their undisclosed attribute
values. The two main security properties are unlinkability, so two presentations
of the same user cannot be tied to one another, and unforgeability, so a user
can only successfully attest to having attribute values that they truly
possess~\cite{EC:CamLys01}.
\begin{outline}
  \item comment: ``define attributes?'' and ``very quick into security properties?'': say what an attribute is (a value about the user that an issuer certifies) before the security properties
  \item then lead more slowly into unlinkability and unforgeability
\end{outline}

Early anonymous credential schemes assume a single issuer certifying all
attributes, creating a single entity responsible for all issuance, and able to
learn all attribute values of every user. Later schemes distribute issuance
across multiple authorities or delegate it to sub-issuers, but the issuer,
though now decentralised, remains conceptually responsible for all attributes.
This reduces trust in any one party, but does not reflect real-world scenarios,
where distinct entities independently certify one or a small subset of
attributes.

In decentralised identity systems, different authorities each typically certify
a small number of attributes: age, residency, citizenship, possession of a
driving licence, possession of a professional qualification. Credential systems
with distributed issuance still require the issuers to coordinate. With many
issuers worldwide, each relevant to only some users, that coordination is
impractical, and between issuers with no relationship to one another it may be
impossible.

\paragraph{Contributions.} In contrast to classical anonymous credential
schemes, we explicitly consider a decentralised setting in which each attribute
is certified by a distinct issuer. We give an aggregate anonymous credential
scheme for this setting: credentials from independent, uncoordinated issuers
combine into one compact credential, provable in one proof. Aggregation is
performed locally by the user without interaction between issuers. The scheme
supports selective disclosure of attribute values.

The key observation behind the scheme is that a credential rests on a chain of
links: from a user's secret key to their public key, from the public key to a
tag, and from the tag to the credentials issued under it. Each link can be
established by a different party. Users show that they know their secret key. A
registration authority, which certifies users' keys in any case, links each
public key to a tag. Issuers then sign attributes under the tag, and their
signatures do not involve the key. At presentation, the verifier checks all three
links together, and so sees that the chain is unbroken without learning the key
or the tag.

The registration authority handles registration only, binding users' keys to
fresh tags. It sees no attribute values and is not involved after registration.
Every credential a user later obtains is issued under their tag. To present
attributes, users aggregate the relevant credentials and prove, in zero
knowledge, that they know the key, hold a valid aggregate credential for the
disclosed values, and that both share a tag. This is a single proof whose cost
grows only with the number of disclosed attributes. Binding by tag stops
different users' credentials being pooled, and lets honest users aggregate
credentials from issuers that have never communicated, without contacting any
issuer again. Separating the requirements on the registration authority from
those on the issuers also lets issuers be stateless. Each user has one tag, so
issuers need not remember whom they have signed for.

The registration authority is not a new party that we add to the model. Every
anonymous credential scheme already assumes that users' keys are certified, so
that credentials cannot be issued to keys that users do not control. That
assumption is usually left implicit, or met by the issuer itself. Here we name
the role.

\subsection*{Related work}

\paragraph{Single-issuer anonymous credential systems.} An early anonymous
credential system issues signatures on committed values, with security under the
strong RSA (sRSA) assumption~\cite{EC:CamLys01,SCN:CamLys02}. A later pairing-based
construction works under the LRSW assumption, though signature size and
verification complexity grow linearly with the number of attribute values
signed~\cite{C:CamLys04}. An extension supports a constrained form of
aggregation, but only in settings where all values are known ahead of
time~\cite{FC:LeeLeeYun13}.
\rebekah{Comments: ``not sure about assumption relevance'' at sRSA and LRSW,
``don't think we should drop assumptions like this, what's the idea? is that
the efficient schemes have bad assumptions'' at $n$-BSDH and $J$-RootDH, and,
on an older printout, ``CDHK15: lots of unnamed assumptions''. The point about
assumptions is now made once, after Comparison. Decide: take the assumption
names out of the related-work prose here and in the CDHK15 and CL11
paragraphs, leaving them in the table only, or keep them and give CDHK15's full
list instead of ``\ldots''.}

Credentials built on the BBS+ signature scheme support selective disclosure,
where a user reveals some of their attribute values and hides the rest, in a
single-issuer setting, with verification linear in the total number of
attributes, and form the basis of the ongoing standardisation effort of the
Internet Engineering Task Force (IETF)~\cite{EC:BonBoy04a,SCN:AuSusMu06}.

Further improvements targeted efficiency. The Pointcheval-Sanders signature
scheme introduces randomisable signatures with sequential aggregation, and
proves unforgeability in the GGM~\cite{RSA:PoiSan16}.
These systems operate in a centralised, single-issuer setting where one
authority certifies all attributes. More recently, structure-preserving
signatures on equivalence classes enable very efficient privacy-preserving
proofs, with security proven in the GGM~\cite{JC:FucHanSla19}. Unlinkable
proofs of knowledge over a subset of signed messages, building on the
Pointcheval-Sanders signature scheme, are also proven secure in the
GGM~\cite{PKC:Sanders20}.

\paragraph{Multi-issuer anonymous credential systems.} Multi-issuer schemes
decentralise the issuer, delegating issuance to sub-issuers or distributing issuance
responsibilities among several
issuers~\cite{TCC:BCKL08,CCS:CamDriDub17,NDSS:SABMD19}. In these cases, the now
decentralised issuer is still responsible for all attributes. While this reduces
the trust placed in a central authority, it does not fully reflect real-world
scenarios where different authorities independently certify distinct attributes.

A multi-issuer credential system supports selective disclosure and
pseudonyms~\cite{AC:CDHK15}. It is proven secure in the universal composability
framework~\cite{FOCS:Canetti01}, under several assumptions including $n$-BSDH
and $J$-RootDH, with a common reference string. Credentials from different
issuers are linked through a pseudonym and not aggregated, so a presentation
contains one proof per credential.

\paragraph{Aggregate anonymous credentials.} Aggregate credentials combine
independently issued credentials into a single compact credential, while
maintaining user privacy during presentations. Indexed aggregate signature
schemes support anonymous credentials this way, with unforgeability under the
$q$-Asymmetric Double Hidden Strong Diffie-Hellman ($q$-ADHSDH) assumption, and
anonymity under the decisional Diffie-Hellman (DDH) assumption in the random
oracle model~\cite{canard2011anonymous}. \rebekah{check CL11 against the paper:
HP22 list it as multi-issuer with a constant-size showing, which contradicts
our table row}

ART-Sign is a traceable aggregate signature scheme using randomisable tags as
pseudonyms, from which H\'ebant and Pointcheval build multi-authority
credentials, the closest scheme to ours~\cite{SCN:HebPoi22}. Showings are of
constant size, and a designated authority can trace users. The signatures are
linearly homomorphic, so two signatures by one issuer under one tag and index
would let anyone sign any value for that index. Issuers sign at most once per
tag and index, and keep a record of what they have signed for each user. In a
bounded variant they instead count the messages signed per tag, and the issuer
key grows with the bound.

Mir et al.\ build issuer-hiding multi-authority credentials from aggregate
signatures with randomisable tags and keys~\cite{CCS:MBGLS23}. Their issuers are
not stateless either. Either the attributes and issuer keys of a user are fixed
when the user's tag is created, or each issuer signs only once per tag. They name the
stateless, dynamic setting as an open question. Unforgeability rests on an
interactive assumption in the random oracle model, or on the GGM, and anonymity
on a new interactive assumption.

\paragraph{Other multi-authority credentials.} Other multi-authority
credentials are in different settings. Threshold issuance without interactive
setup lets uncoordinated issuers jointly certify one
credential~\cite{EPRINT:GGPR25}\rebekah{``don't cite eprint'': cryptobib has
no published version of GGPR25 (only ePrint 2025/2042); check whether it has
appeared since, else keep the ePrint key}, while here each issuer certifies its own
attributes. Independent authorities can also each sign a user's global
identity, with one proof per credential~\cite{anada2019decentralized}.
Credentials can be derived from existing identity documents with zk-SNARKs,
without new issuer infrastructure, at the cost of a general-purpose proof per
credential~\cite{SP:RWGM23}. Issuer hiding, where a verifier learns only that an
issuer lies in a set it accepts, is an orthogonal
problem~\cite{CANS:BEKRS21,PoPETS:SanTra24}, and one we would like to address in
future work.

\paragraph{Comparison.} H\'ebant and Pointcheval's certification authority plays
the same part as our registration authority, binding each user to one tag at
registration. The difference is what the tag is. In ART-Sign the tag is the
user's public key, whose validity is proven again at every showing. Ours is a
separate value that the registration authority signs together with the key, so
issuers keep no record, and presentations prove knowledge of the registration
signature instead.
ART-Sign offers traceability, which our scheme does not.
\rebekah{Comments on the Comparison bullets: ``really compress, 3 sentences in
total?'', and on the bullets after it ``last bullet is useful, others should be
mostly covered by comparison higher up. just one sentence here about ART, one
about the others.'' Cut this paragraph to one sentence on ART-Sign, and the next
paragraph to one sentence on the other schemes, keeping the assumptions point.}

\Cref{tab:agg-creds-related} compares our scheme with several single- and
multi-issuer credential systems, reporting asymptotic complexities. In the
single-issuer schemes the question of issuer state does not arise, and the
issuers of CDHK15 keep no per-user record either, but its credentials are not
aggregated. The efficient schemes in the table rely on strong assumptions:
$q$-type assumptions, interactive assumptions, or the GGM. For SPS schemes with
short signatures this cannot be avoided, since they cannot be proven secure from
non-interactive assumptions by a large class of reductions~\cite{AC:AbeGroOhk11}.
Our tag-based signature scheme needs only BCDH, and the GGM enters only through
the registration signature.

\begin{table}[htbp] \centering \small
  \begin{tabular}{l|c|c|c|c|c}
    \toprule Scheme & Assumptions & Multi-issuer & $|$cred$|$ & Prover cost &
Issuer state \\ \midrule
    CL02~\cite{SCN:CamLys02} & sRSA & No & $O(1)$ & $O(n)$ & -- \\
    CL04~\cite{C:CamLys04} & LRSW & No & $O(n)$ & $O(n)$ & -- \\
    BBS+~\cite{EC:BonBoy04a,SCN:AuSusMu06} & $q$-SDH & No & $O(1)$ & $O(n)$ & -- \\
    CL11~\cite{canard2011anonymous} & $q$-ADHSDH & Partial & $O(n)$ & $O(k)$ & ? \\
    CDHK15~\cite{AC:CDHK15} & $n$-BSDH, $J$-RootDH, \ldots & Yes & $O(K)$ & $O(K)$ & None \\
    HP22~\cite{SCN:HebPoi22} & GGM & Yes & $O(1)$ & $O(k)$ & Per tag \\
    MBGLS23~\cite{CCS:MBGLS23} & Interactive or GGM & Yes & $O(1)$ & $O(k)$ & Per tag \\
    Our scheme & BCDH, DDH, ROM, GGM & Yes & $O(1)$ & $O(k)$ & None \\
\bottomrule
  \end{tabular}
  \caption[Comparison with existing anonymous credential schemes]{Comparison with
existing single and multi-issuer anonymous credential schemes. $n$ is the total
number of attributes a scheme supports, $k$ the number of attributes disclosed in
a presentation, $K$ the number of issuers in a presentation, and $|$cred$|$ the
size of the credentials needed for a presentation. The assumptions of the other
schemes are those of the cited papers.
Issuer state is what an issuer has to remember about each user. For HP22 this
is which tags it has signed for, and for MBGLS23 either the same or the
attributes and issuer keys fixed in advance.} \label{tab:agg-creds-related} \end{table}

\paragraph{Layout.} The rest of the chapter is organised as follows. \Cref{sec:agg-creds-definitions} states the problem and
formalises it as an ideal functionality.  \Cref{sec:agg-creds-construction}
gives an abstract construction from generic primitives.
\Cref{sec:agg-creds-instantiation} gives a concrete instantiation, including our
tag-based aggregate signature scheme, which \Cref{sec:agg-creds-security} proves
secure.  \Cref{sec:agg-creds-evaluation} evaluates performance, and
\Cref{sec:agg-creds-summary} summarises the chapter.
